\documentclass[12pt]{article}
\usepackage{graphicx,amsmath,amsfonts,amssymb,epsfig,euscript,enumerate}
\usepackage[T1]{fontenc}
\usepackage[utf8x]{inputenc}

\newtheorem{exercise}{Ejercicio}
\newcommand{\bej}{\begin{exercise}\rm}
\newcommand{\fej}{\end{exercise}}

\newcommand{\R}{\mathbb{R}}
\newcommand{\C}{\mathbb{C}}
\def\dt{\Delta t}
\def\dx{\Delta x}

\topmargin-2cm \vsize 29.5cm \hsize 21cm
\setlength{\textwidth}{16.75cm}\setlength{\textheight}{23.5cm}
\setlength{\oddsidemargin}{0.0cm}
\setlength{\evensidemargin}{0.0cm}

\begin{document}
\centerline{{\small Universidad de Buenos Aires - Facultad de Ciencias Exactas y Naturales - Depto. de Matemática}}
 
 \vskip 0.2cm
 \hrulefill
 \vskip 0.2cm

 \centerline{{\bf\Huge {\sc Elementos de Cálculo Numérico}}}
 \vskip 0.2cm
 \centerline{\ttfamily Primer Cuatrimestre 2026}
 \hrulefill

 \bigskip
 \centerline{\bf Práctica N$^\circ$ 4 (segunda parte): Polinomios ortogonales}
 \bigskip

\section*{Polinomios de Tchebyshev}

\subsection*{Aproximación de funciones}

\bej
(Ortogonalidad y coeficientes de Chebyshev)
Los polinomios de Chebyshev de primera especie se definen por $T_n(x)=\cos(n\arccos x)$ en $[-1,1]$.
\begin{enumerate}[(a)]
    \item Demuestre que son ortogonales con el peso $(1-x^2)^{-1/2}$:
    \[
    \int_{-1}^{1}\frac{T_m(x)T_n(x)}{\sqrt{1-x^2}}\,dx=0\quad (m\neq n).
    \]
    \item Muestre que los coeficientes de Chebyshev de $f$ coinciden con coeficientes de Fourier de $f(\cos\theta)$.
\end{enumerate}
\fej

\bej
(Interpolación de Chebyshev y estabilidad)
\begin{enumerate}[(a)]
    \item Defina la constante de Lebesgue para interpolación polinomial y compare cualitativamente nodos equiespaciados vs nodos de Chebyshev.
    \item Explique por qué la elección de nodos de Chebyshev evita el fenómeno de Runge.
\end{enumerate}
\fej

\subsection*{II.2 Métodos pseudo-espectrales}

\subsection*{Método espectral de Chebyshev para ODEs con borde}

\bej
(Problema modelo con Dirichlet: orden de truncamiento)
Considere
\[
-u''(x)+u(x)=f(x),\qquad x\in[-1,1],
\qquad
u(-1)=u(1)=0.
\]
Suponga que se aproxima con un método espectral de Chebyshev (tau/colocación) usando $N+1$ modos.
\begin{enumerate}[(a)]
    \item Escriba la expansión truncada
    \[
    u_N(x)=\sum_{n=0}^{N} a_n T_n(x),
    \qquad
    f_N(x)=\sum_{n=0}^{N} b_n T_n(x),
    \]
    e indique cómo se imponen las condiciones de borde en el sistema lineal.
    \item Sabiendo (por el apunte) que si $f\in C^r([-1,1])$ entonces sus coeficientes de Chebyshev decaen con orden $r$, estime el orden esperado del error de truncamiento de $u_N$.
    \item Compare el orden de convergencia de la aproximación de $u$ con el de la aproximación de $f$ y discuta si hay ganancia, pérdida o conservación de orden.
\end{enumerate}
\fej

\subsection*{II.3 Cuadratura en nodos de Chebyshev}

\bej
(Mismo esquema con condiciones de Neumann)
Considere
\[
-u''(x)+u(x)=f(x),\qquad x\in[-1,1],
\qquad
u'(-1)=u'(1)=0.
\]
\begin{enumerate}[(a)]
    \item Proponga la formulación espectral de Chebyshev con $N+1$ modos e indique cómo se incorporan las condiciones de Neumann en el sistema.
    \item Explique qué cambia respecto del caso Dirichlet en términos de estabilidad numérica y condicionamiento al crecer $N$.
    \item Bajo la hipótesis $f\in C^r([-1,1])$, compare el orden esperado de convergencia con el del ejercicio anterior.
\end{enumerate}
\fej

\bej
(Adaptación del método a otras condiciones de borde)
Considere
\[
\mathcal L u = f \quad \text{en }[-1,1],
\]
donde $\mathcal L$ es un operador diferencial lineal de segundo orden.
\begin{enumerate}[(a)]
    \item Para condiciones de Robin
    \[
    \alpha_-u(-1)+\beta_-u'(-1)=g_-,
    \qquad
    \alpha_+u(1)+\beta_+u'(1)=g_+,
    \]
    describa dos estrategias para adaptar el método espectral de Chebyshev:
    \begin{enumerate}[i.]
        \item imponer las condiciones en el sistema (tau/colocación),
        \item construir una base que ya satisfaga las condiciones de borde.
    \end{enumerate}
    \item Elija una de las estrategias y escriba explícitamente el sistema lineal para los coeficientes truncados.
    \item Discuta qué parte del análisis de orden de truncamiento se mantiene igual y qué parte depende del tipo de borde.
\end{enumerate}
\fej

\bej
(Nodos de Chebyshev-Gauss-Lobatto y matriz de diferenciación)
Considere los nodos
\[
x_j=\cos\left(\frac{\pi j}{N}\right),\qquad j=0,\ldots,N.
\]
Sea
\[
\lambda_j=\frac{1}{\prod_{m\ne j}(x_j-x_m)}
\]
los pesos baricéntricos asociados a esos nodos.
Sea $p_N\in\mathcal P_N$ el interpolante en estos nodos y $v_j=p_N(x_j)$.
\begin{enumerate}[(a)]
    \item Muestre que existe una matriz $D\in\mathbb R^{(N+1)\times(N+1)}$ tal que
    \[
    p_N'(x_i)=\sum_{j=0}^{N}D_{ij}v_j,
    \qquad i=0,\ldots,N.
    \]
    \item Deduzca, en términos de los pesos baricéntricos, que
    \[
    D_{ij}=\frac{\lambda_j}{\lambda_i}\frac{1}{x_i-x_j}\quad(i\ne j),
    \]
    y exprese los términos diagonales mediante
    \[
    D_{ii}=-\sum_{j\ne i}D_{ij}.
    \]
    (Opcional: recupere las fórmulas cerradas en nodos de Chebyshev-Gauss-Lobatto.)
    \item Verifique que $D\mathbf 1=0$ y discuta por qué esto es consistente con la derivada de una constante.
\end{enumerate}
\fej

\bej
(Matriz de segunda derivada y reescalado de intervalo)
Sea $D$ la matriz de diferenciación del ejercicio anterior y $D^{(2)}:=D^2$.
\begin{enumerate}[(a)]
    \item Explique por qué, para una aproximación pseudo-espectral en nodos de Chebyshev, se puede aproximar
    \[
    u''(x_i)\approx \sum_{j=0}^{N}D^{(2)}_{ij}u(x_j).
    \]
    \item Si el problema está en $[a,b]$ con cambio afín $x=\frac{2t-(a+b)}{b-a}$, deduzca cómo escalan las matrices de derivación para $\frac{d}{dt}$ y $\frac{d^2}{dt^2}$.
    \item Aplique lo anterior para escribir el sistema lineal de colocación de
    \[
    -u''(t)+u(t)=f(t),\qquad t\in[a,b],
    \]
    con condiciones de Dirichlet en los extremos.
\end{enumerate}
\fej

\bej
(Transformada de Chebyshev vía FFT/DCT)
Sea $x_j=\cos(\pi j/N)$ y $v_j=u(x_j)$.
\begin{enumerate}[(a)]
    \item Muestre que el cálculo de coeficientes de Chebyshev de $u_N(x)=\sum_{n=0}^{N}a_nT_n(x)$ puede reducirse a una DFT real (equivalentemente, a una DCT), usando la variable angular $x=\cos\theta$.
    \item Describa un algoritmo $O(N\log N)$ para pasar de valores nodales $\{v_j\}$ a coeficientes $\{a_n\}$ y viceversa.
    \item Escriba cómo obtener los coeficientes de $u_N'$ a partir de $\{a_n\}$ (recurrencia espectral) y deduzca un algoritmo pseudo-espectral para computar $u'(x_j)$ en costo $O(N\log N)$.
\end{enumerate}
\fej

\bej
(Interpolación baricéntrica de Chebyshev en costo $O(N)$ por evaluación)
Considere los nodos de Chebyshev-Gauss-Lobatto
\[
x_j=\cos\left(\frac{\pi j}{N}\right),\qquad j=0,\ldots,N,
\]
y los pesos baricéntricos explícitos
\[
\lambda_j=\frac{1}{\prod_{m\ne j}(x_j-x_m)}
=\frac{2^{N-1}}{N}(-1)^j\delta_j,
\qquad
\delta_0=\delta_N=\frac12,\ \delta_j=1\ (1\le j\le N-1).
\]
Sea $p_N$ el interpolante de datos $f_j=f(x_j)$.
\begin{enumerate}[(a)]
    \item Escriba la fórmula baricéntrica de $p_N(x)$ y el tratamiento del caso $x=x_j$.
    \item Muestre que, dados $\{f_j\}$ y $\{\lambda_j\}$, evaluar $p_N(x)$ en un punto arbitrario $x\in[-1,1]$ cuesta $O(N)$ operaciones.
    \item Si se quieren evaluar $M$ puntos, estime el costo total y compare con resolver el sistema de Vandermonde de interpolación.
    \item \textbf{Pregunta conceptual}: ¿es lo mismo evaluar/interpolar con la fórmula baricéntrica que calcular los coeficientes de la expansión de Chebyshev
    \[
    p_N(x)=\sum_{n=0}^{N}a_nT_n(x)?
    \]
    Justifique brevemente.
\end{enumerate}
\fej

\bej
(Cuadratura de Fejér a partir de pesos baricéntricos)
Sea $x_j=\cos\big((2j+1)\pi/(2N)\big)$, $j=0,\ldots,N-1$ (nodos de Fejér de primera especie), y $\lambda_j$ sus pesos baricéntricos.
\begin{enumerate}[(a)]
    \item Recuerde que la cuadratura interpolatoria asociada tiene pesos
    \[
    w_j=\int_{-1}^{1}\ell_j(x)\,dx,
    \]
    donde $\ell_j$ son los polinomios cardinales de Lagrange. Escriba $\ell_j$ en forma baricéntrica usando $\lambda_j$.
    \item Explique cómo obtener los pesos $w_j$ a partir de los valores de $\ell_j$ en los nodos y una transformada coseno (DCT/FFT), en lugar de integrar simbólicamente cada $\ell_j$.
    \item Compare costos: método directo $O(N^2)$ vs método basado en transformadas rápidas $O(N\log N)$.
    \item Verifique en qué grado de polinomios la regla resultante es exacta (grado de precisión de Fejér I).
\end{enumerate}
\fej

\bej
(Cuadratura de Gauss-Chebyshev: nodos abiertos)
Considere la integral con peso
\[
I(f)=\int_{-1}^{1}\frac{f(x)}{\sqrt{1-x^2}}\,dx.
\]
\begin{enumerate}[(a)]
    \item Muestre que, usando los nodos abiertos
    \[
    x_j=\cos\left(\frac{2j+1}{2N}\pi\right),\qquad j=0,\ldots,N-1,
    \]
    se obtiene la regla
    \[
    I(f)\approx Q_N^{\mathrm{GC}}(f):=\frac{\pi}{N}\sum_{j=0}^{N-1}f(x_j).
    \]
    \item Verifique el grado de exactitud de la regla (en el espacio de polinomios con el peso dado).
    \item Aplique la regla para aproximar
    \[
    \int_{-1}^{1}\frac{e^x}{\sqrt{1-x^2}}\,dx,
    \]
    y discuta la convergencia al aumentar $N$.
\end{enumerate}
\fej

\bej
(Cuadratura de Clenshaw-Curtis: nodos cerrados)
Considere la integral sin peso
\[
J(f)=\int_{-1}^{1}f(x)\,dx,
\]
y los nodos cerrados de Chebyshev-Gauss-Lobatto
\[
x_j=\cos\left(\frac{\pi j}{N}\right),\qquad j=0,\ldots,N.
\]
\begin{enumerate}[(a)]
    \item Defina la regla de Clenshaw-Curtis
    \[
    J(f)\approx Q_N^{\mathrm{CC}}(f)=\sum_{j=0}^{N}w_j f(x_j),
    \]
    con $w_j=\int_{-1}^{1}\ell_j(x)\,dx$.
    \item Explique cómo calcular los pesos $\{w_j\}$ eficientemente a partir de una DCT/FFT (sin integrar cada $\ell_j$ en forma simbólica).
    \item Compare el costo de preprocesamiento de pesos con un método directo y luego el costo por evaluación de la cuadratura.
    \item Compare cualitativamente esta regla (nodos cerrados) con Gauss-Chebyshev/Fejér (nodos abiertos): estabilidad cerca de extremos y facilidad para imponer condiciones de borde en métodos pseudo-espectrales.
\end{enumerate}
\fej

\bej
(Relación entre Fejér, Clenshaw-Curtis y transformadas rápidas)
\begin{enumerate}[(a)]
    \item Muestre que tanto Fejér como Clenshaw-Curtis pueden interpretarse como: interpolar en nodos de Chebyshev y luego integrar el interpolante.
    \item Explique por qué ambas reglas admiten implementación $O(N\log N)$ usando FFT/DCT.
    \item Discuta en qué situaciones conviene usar nodos cerrados (Clenshaw-Curtis) o abiertos (Fejér/Gauss-Chebyshev).
\end{enumerate}
\fej

\bej
(Cálculo de nodos de cuadratura en familias de Chebyshev: costo $O(N^2)$)
Para cada familia de nodos de Chebyshev (Gauss-Chebyshev, Clenshaw-Curtis/Gauss-Lobatto, Fejér I y Fejér II):
\begin{enumerate}[(a)]
    \item Escriba una fórmula explícita para los nodos en términos de cosenos.
    \item Proponga un algoritmo que calcule \emph{todos} los nodos usando esa fórmula con costo total $O(N^2)$ (por ejemplo, evaluando trigonometría con un bucle doble o mediante recurrencias por nodo).
    \item Compare con una implementación directa de las fórmulas explícitas en costo $O(N)$ y discuta por qué, aun así, el costo dominante de la cuadratura puede estar en el cálculo de pesos o en la evaluación de la función.
    \item (Opcional) Discuta cómo mejorar el cálculo de nodos y pesos cuando se requiere una secuencia de tamaños $N=2^m$.
\end{enumerate}
\fej

\bej
(Nodos anidados de Clenshaw-Curtis y estimación de error embebida)
Considere Clenshaw-Curtis con
\[
x_j^{(N)}=\cos\left(\frac{\pi j}{N}\right),\qquad j=0,\ldots,N.
\]
\begin{enumerate}[(a)]
    \item Muestre que, si se pasa de $N$ a $2N$, se cumple
    \[
    x_j^{(N)}=x_{2j}^{(2N)},\qquad j=0,\ldots,N,
    \]
    es decir, los nodos son anidados.
    \item Explique cómo reutilizar evaluaciones de $f$ al pasar de $Q_N^{\mathrm{CC}}(f)$ a $Q_{2N}^{\mathrm{CC}}(f)$.
    \item Proponga un estimador de error embebido del tipo
    \[
    E_N:=\left|Q_{2N}^{\mathrm{CC}}(f)-Q_N^{\mathrm{CC}}(f)\right|,
    \]
    y describa un criterio adaptativo de parada para alcanzar una tolerancia dada.
    \item Discuta ventajas y limitaciones de este estimador cuando $f$ tiene baja regularidad o singularidades cerca de los extremos.
\end{enumerate}
\fej

\section*{III. Polinomios de Legendre y cuadratura gaussiana}

\bej
(Mejor aproximación en $L^2$)
Consideremos el espacio de funciones de cuadrado integrable en $[-1, 1]$ con el producto interno estándar ($w(x)=1$). Queremos encontrar el polinomio $p(x)$ de grado a lo sumo 1 que mejor aproxima a la función $f(x) = e^x$ en el sentido de los cuadrados mínimos (norma $L^2$).
\begin{enumerate}[(a)]
    \item Recuerde que los primeros polinomios de Legendre son $P_0(x) = 1$ y $P_1(x) = x$. Verifique que son ortogonales en $[-1, 1]$.
    \item Calcule las normas $\|P_0\|^2$ y $\|P_1\|^2$.
    \item La proyección ortogonal de $f$ sobre el subespacio generado por $\{P_0, P_1\}$ está dada por:
    \[
    p(x) = \frac{\langle f, P_0 \rangle}{\|P_0\|^2} P_0(x) + \frac{\langle f, P_1 \rangle}{\|P_1\|^2} P_1(x).
    \]
    Calcule explícitamente este polinomio para $f(x) = e^x$.
    \item Compare cualitativamente este polinomio con:
    \begin{itemize}
        \item El polinomio de Taylor de grado 1 centrado en $x=0$.
        \item El polinomio interpolador de Lagrange en los nodos $-1$ y $1$.
    \end{itemize}
\end{enumerate}
\fej

\bej
(Construcción de Polinomios Ortogonales)
Se desea construir una familia de polinomios $\{q_0, q_1, q_2\}$ ortogonales en el intervalo $[0, 1]$ respecto al peso $w(x) = 1$.
\begin{enumerate}
    \item Utilice el proceso de Gram-Schmidt comenzando con la base canónica $\{1, x, x^2\}$ para obtener $q_0, q_1, q_2$.
    \item Verifique que las raíces de $q_2(x)$ son reales, distintas y están dentro del intervalo $(0, 1)$.
\end{enumerate}
\fej

\bej
(Cuadratura Gaussiana)
Queremos aproximar la integral $I(f) = \int_{-1}^1 f(x)\,dx$ mediante una regla de cuadratura de 2 puntos:
\[
Q(f) = w_1 f(x_1) + w_2 f(x_2).
\]
\begin{enumerate}
    \item Sabemos que la cuadratura de Gauss con $n$ nodos es exacta para polinomios de grado $2n-1$. Para $n=2$, esto significa que debe ser exacta para polinomios hasta grado 3.
    \item Utilice las raíces del polinomio de Legendre $P_2(x)$ (hallado/deducido en ejercicios anteriores) como nodos $x_1, x_2$.
    \item Determine los pesos $w_1, w_2$ imponiendo que la regla sea exacta para $f(x)=1$ y $f(x)=x$ (o usando la fórmula de interpolación de Lagrange para los pesos).
    \item Verifique que la regla obtenida integra exactamente a $x^2$ y $x^3$.
    \item Use esta regla para aproximar $\int_{-1}^1 e^x dx$ y compare con el valor exacto.
\end{enumerate}
\fej

\bej
(Método de Golub--Welsch para cuadratura de Gauss-Legendre)
Sea $\{P_n\}$ la familia de polinomios de Legendre ortogonales en $[-1,1]$ y considere la recurrencia tritérmino normalizada.
\begin{enumerate}[(a)]
    \item Construya la matriz de Jacobi $J_n\in\mathbb R^{n\times n}$ asociada a dicha recurrencia.
    \item Muestre que los nodos de Gauss-Legendre de orden $n$ son los autovalores de $J_n$.
    \item Muestre que los pesos pueden obtenerse a partir de la primera componente de los autovectores normalizados de $J_n$.
    \item Describa el algoritmo de Golub--Welsch y estime su costo computacional.
\end{enumerate}
\fej


\bej
(Interpolación vs Proyección)
Los coeficientes de la serie de Chebyshev de $f$ están dados por las integrales $c_k = \frac{2}{\pi} \int_{-1}^1 f(x) T_k(x) (1-x^2)^{-1/2} dx$ (con $c_0$ dividido por 2).
Considere la fórmula de cuadratura de Gauss-Chebyshev con $M$ puntos, que es exacta para polinomios de grado $2M-1$:
\[ \int_{-1}^1 \frac{g(x)}{\sqrt{1-x^2}} dx \approx \frac{\pi}{M} \sum_{j=1}^M g(x_j), \]
donde $x_j$ son las raíces de $T_M(x)$.
\begin{enumerate}[(a)]
    \item Aproxime la integral de los coeficientes $c_k$ usando esta regla de cuadratura con $M=N+1$ puntos.
    \item Muestre que los coeficientes aproximados $\tilde{c}_k$ coinciden con los coeficientes del polinomio interpolador de $f$ en los nodos de Chebyshev $x_0, \dots, x_N$, expresado en la base $\{T_0, \dots, T_N\}$.
    \item Concluya que calcular la interpolación en nodos de Chebyshev es equivalente a discretizar la proyección ortogonal mediante cuadratura Gaussiana.
\end{enumerate}
\fej

\begin{exercise}[Constante de Lebesgue]
La constante de Lebesgue $\Lambda_n$ mide qué tan bien condicionado está el problema de interpolación.
\begin{enumerate}[(a)]
\item Defina la constante de Lebesgue:
\[
\Lambda_n = \max_{x \in [a,b]} \sum_{k=0}^n |\ell_k(x)|.
\]
\item Demuestre que el error de interpolación satisface:
\[
\|f - p_n\|_\infty \leq (1 + \Lambda_n) \|f - p_n^*\|_\infty,
\]
donde $p_n^*$ es el mejor polinomio de aproximación uniforme de grado $n$.
\item Investigue y enuncie los valores de $\Lambda_n$ para:
   \begin{itemize}
   \item Nodos equiespaciados: $\Lambda_n = O(2^n / n)$ (crece exponencialmente).
   \item Nodos de Chebyshev: $\Lambda_n = O(\log n)$ (crece logarítmicamente).
   \end{itemize}
\item Explique por qué una constante de Lebesgue grande indica que la interpolación puede amplificar errores en los datos.
\end{enumerate}
\end{exercise}

\begin{exercise}[Fenómeno de Runge]
El fenómeno de Runge muestra que aumentar el grado de interpolación con nodos equiespaciados puede hacer diverger el polinomio interpolante.
\begin{enumerate}[(a)]
\item Considere la función de Runge $f(x) = \frac{1}{1 + 25x^2}$ en $[-1, 1]$.
\item Explique por qué esta función, siendo analítica en $\mathbb{R}$, presenta problemas de convergencia en interpolación polinomial.
\textbf{Sugerencia}: Considere los polos complejos $x = \pm i/5$ y su proximidad al intervalo real.
\item Para nodos equiespaciados $x_k = -1 + 2k/n$ en $[-1, 1]$, el producto:
\[
\prod_{k=0}^n (x - x_k)
\]
crece exponencialmente cuando $x$ se acerca a los extremos $\pm 1$. Estime el orden de magnitud de este producto cerca de $x = 1$.
\item Usando el teorema del error de interpolación, explique por qué aunque $f$ sea suave, el error de interpolación puede crecer con $n$ para nodos equiespaciados.
\item ¿Por qué el fenómeno de Runge no ocurre con nodos de Chebyshev?
\end{enumerate}
\end{exercise}

\begin{exercise}[Nodos de Chebyshev]
Los nodos de Chebyshev resuelven el fenómeno de Runge al minimizar el error de interpolación.
\begin{enumerate}[(a)]
\item Los nodos de Chebyshev en $[-1, 1]$ se definen como:
\[
x_k = \cos\left(\frac{(2k+1)\pi}{2(n+1)}\right), \quad k = 0, \ldots, n.
\]
Calcule los nodos de Chebyshev para $n = 4$ (5 nodos).
\item Verifique que estos nodos son las raíces del polinomio de Chebyshev $T_{n+1}(x)$.
\item Observe que los nodos se acumulan cerca de los extremos $x = \pm 1$ y están más espaciados en el centro. ¿Por qué esta distribución es óptima?
\item Demuestre que con nodos de Chebyshev:
\[
\left|\prod_{k=0}^n (x - x_k)\right| \leq \frac{1}{2^n}, \quad \forall x \in [-1, 1].
\]
\textbf{Sugerencia}: Use que $\prod_{k=0}^n (x - x_k) = \frac{1}{2^n} T_{n+1}(x)$ y que $|T_{n+1}(x)| \leq 1$ en $[-1,1]$.
\item Explique por qué esta cota garantiza convergencia para funciones suficientemente suaves, a diferencia de nodos equiespaciados.
\end{enumerate}
\end{exercise}

\begin{exercise}[Error de interpolación con nodos de Chebyshev]
\begin{enumerate}[(a)]
\item Enuncie el teorema: Si $f \in C^{n+1}[-1,1]$ y $p_n$ interpola $f$ en los nodos de Chebyshev, entonces:
\[
\|f - p_n\|_\infty \leq \frac{\|f^{(n+1)}\|_\infty}{2^n (n+1)!}.
\]
\item Compare esta cota con la de nodos equiespaciados donde $\prod_{k=0}^n (x - x_k)$ puede crecer exponencialmente.
\item Para $f(x) = \sin(x)$ en $[-1, 1]$, estime el error máximo de interpolación con $n = 10$ nodos de Chebyshev.
\item ¿Para qué valor de $n$ el error es menor que $10^{-10}$?
\item Investigue sobre convergencia para funciones analíticas: la interpolación de Chebyshev converge geométricamente (exponencialmente) para funciones analíticas.
\end{enumerate}
\end{exercise}


\end{document}
